\chapter{Configuración de la campaña principal}
\label{anexo:campania}

La configuración completa se encuentra en el directorio \path{campaigns}, en el fichero \path{/hpc-tfg-main-campaign.json}. Los parámetros determinantes para el análisis son:

\begin{lstlisting}
{
  "campaign_name": "hpc_tfg_main_campaign",
  "launcher": "slurm",
  "partition": "guest",
  "nodelist": "compute-0-4",
  "duration": 120,
  "monitoring": {
    "enabled": true,
    "publish": {
      "internal_metrics": true,
      "external_metrics": true
    },
    "grouping": {"mode": "run_id"}
  },
  "external_measurement": {
    "enabled": true,
    "provider": "vampire"
  },
  "languages": ["c"],
  "benchmarks": [
    "idle",
    "memory-stream-hpc",
    "compute-dgemm-hpc"
  ],
  "reps": [1, 2, 3, 4, 5]
}
\end{lstlisting}

La configuración completa mantiene fuera de este extracto la URL del servicio y las credenciales. En la campaña final, \texttt{compute-0-4} se asocia con el dispositivo Vampire \texttt{hpm4}; la captura reserva 2~s antes y 15~s después del benchmark y reintenta la descarga si las muestras todavía no cubren la ventana completa. Las métricas internas y externas se publican con agrupación por \texttt{run\_id}.

Los perfiles \texttt{p1}, \texttt{p2}, \texttt{p4}, \texttt{p6}, \texttt{p8} y \texttt{p12} mantienen un hilo por proceso y solicitan, respectivamente, 1, 2, 4, 6, 8 y 12 procesos y CPU. Los tamaños utilizados son 50 millones de elementos por vector en STREAM y matrices de orden 1024 en DGEMM.
